\subsection*{Cluster Values of Sequences}
\begin{exposition}
A convergent sequence has only one limiting value. A bounded divergent
sequence may still have values that it approaches along selected subsequences.
Cluster values name this repeated limiting behavior without requiring the
entire sequence to converge.
\end{exposition}

% ---------------------------------------------------------
% Definition --- Cluster Value of a Sequence
% ---------------------------------------------------------

\begin{definitionbox}{Definition (Cluster Value of a Sequence)}
\begin{definition}[Cluster Value of a Sequence]
\label{def:cluster-value-sequence}
Let $(x_n)$ be a real sequence, and let $L\in\mathbb{R}$. The number $L$ is
a \textbf{cluster value} of $(x_n)$ if for every $\varepsilon>0$ and every
$N\in\mathbb{N}$ there exists $n\ge N$ such that
\[
  |x_n-L|<\varepsilon.
\]
\end{definition}
\end{definitionbox}

\begin{remark*}[Standard quantified statement]
\[
L\text{ is a cluster value of }(x_n)
\Longleftrightarrow
\forall \varepsilon>0\;\forall N\in\mathbb{N}\;\exists n\ge N
\;\bigl(|x_n-L|<\varepsilon\bigr).
\]
\end{remark*}

\begin{remark*}[Definition predicate reading]
\[
\operatorname{SequenceClusterValue}(x_n,L)
\coloneqq
\forall \varepsilon>0\;\forall N\in\mathbb{N}\;\exists n\ge N
\;\bigl(|x_n-L|<\varepsilon\bigr).
\]
\end{remark*}

\begin{remark*}[Negated quantified statement]
\[
L\text{ is not a cluster value of }(x_n)
\Longleftrightarrow
\exists \varepsilon>0\;\exists N\in\mathbb{N}\;\forall n\ge N
\;\bigl(|x_n-L|\ge\varepsilon\bigr).
\]
\end{remark*}

\begin{remark*}[Interpretation]
No matter how far out in the sequence one looks, some later term returns
inside every neighborhood of $L$.
\end{remark*}

\begin{remark*}[Dependencies]\hfill
\begin{itemize}
  \item \hyperref[def:sequence]{Sequence}
\end{itemize}
\end{remark*}

% ---------------------------------------------------------
% Theorem --- Cluster Values Are Subsequential Limits
% ---------------------------------------------------------

\begin{theorembox}{Theorem (Cluster Values Are Subsequential Limits)}
\begin{theorem}[Cluster Values Are Subsequential Limits]
\label{thm:cluster-values-are-subsequential-limits}
Let $(x_n)$ be a real sequence, and let $L\in\mathbb{R}$. Then $L$ is a
cluster value of $(x_n)$ if and only if $L$ is a subsequential limit of
$(x_n)$.

\smallskip
\noindent
\hyperref[prf:cluster-values-are-subsequential-limits]{\textit{Go to proof.}}
\end{theorem}
\end{theorembox}

\begin{remark*}[Standard quantified statement]
\[
\forall (x_n)\subseteq\mathbb{R}\;\forall L\in\mathbb{R}\;
\Bigl(
L\text{ is a cluster value of }(x_n)
\Longleftrightarrow
L\text{ is a subsequential limit of }(x_n)
\Bigr).
\]
\end{remark*}

\begin{remark*}[Theorem predicate reading]
\[
\operatorname{SequenceClusterValue}(x_n,L)
\Longleftrightarrow
\operatorname{SubsequentialLimit}(x_n,L).
\]
\end{remark*}

\begin{remark*}[Interpretation]
The neighborhood formulation and the subsequence formulation describe the
same phenomenon: infinitely many opportunities to get arbitrarily close to
$L$.
\end{remark*}

\begin{remark*}[Dependencies]\hfill
\begin{itemize}
  \item \hyperref[def:cluster-value-sequence]{Cluster Value of a Sequence}
  \item \hyperref[def:subsequential-limit]{Subsequential Limit}
  \item \hyperref[thm:frequent-properties-yield-subsequences]{Frequent Properties Yield Subsequences}
\end{itemize}
\end{remark*}

% ---------------------------------------------------------
% Theorem --- Bounded Sequences Have Cluster Values
% ---------------------------------------------------------

\begin{theorembox}{Theorem (Bounded Sequences Have Cluster Values)}
\begin{theorem}[Bounded Sequences Have Cluster Values]
\label{thm:bounded-sequences-have-cluster-values}
Every bounded real sequence has at least one cluster value.

\smallskip
\noindent
\hyperref[prf:bounded-sequences-have-cluster-values]{\textit{Go to proof.}}
\end{theorem}
\end{theorembox}

\begin{remark*}[Standard quantified statement]
\[
\forall (x_n)\subseteq\mathbb{R}\;
\Bigl(
  \bigl[\exists M>0\;\forall n\in\mathbb{N}\;(|x_n|\le M)\bigr]
  \Longrightarrow
  \exists L\in\mathbb{R}\;
  \bigl(L\text{ is a cluster value of }(x_n)\bigr)
\Bigr).
\]
\end{remark*}

\begin{remark*}[Theorem predicate reading]
\[
\operatorname{BoundedSequence}(x_n)
\Longrightarrow
\exists L\in\mathbb{R}\;\operatorname{SequenceClusterValue}(x_n,L).
\]
\end{remark*}

\begin{remark*}[Interpretation]
Boundedness traps the sequence in a finite interval, and
Bolzano-Weierstrass extracts a convergent subsequence from that trap.
\end{remark*}

\begin{remark*}[Dependencies]\hfill
\begin{itemize}
  \item \hyperref[def:bounded-sequence]{Bounded Sequence}
  \item \hyperref[thm:bolzano-weierstrass-sequences]{Bolzano-Weierstrass Theorem for Sequences}
  \item \hyperref[thm:cluster-values-are-subsequential-limits]{Cluster Values Are Subsequential Limits}
\end{itemize}
\end{remark*}

% ---------------------------------------------------------
% Theorem --- Limsup and Liminf Are Extremal Cluster Values
% ---------------------------------------------------------

\begin{theorembox}{Theorem (Limsup and Liminf Are Extremal Cluster Values)}
\begin{theorem}[Limsup and Liminf Are Extremal Cluster Values]
\label{thm:limsup-liminf-extremal-cluster-values}
Let $(x_n)$ be a bounded real sequence. Then
$\limsup_{n\to\infty}x_n$ is the largest cluster value of $(x_n)$, and
$\liminf_{n\to\infty}x_n$ is the smallest cluster value of $(x_n)$.

\smallskip
\noindent
\hyperref[prf:limsup-liminf-extremal-cluster-values]{\textit{Go to proof.}}
\end{theorem}
\end{theorembox}

\begin{remark*}[Standard quantified statement]
\[
\begin{aligned}
&\forall (x_n)\subseteq\mathbb{R}\;
\Bigl(
  (x_n)\text{ is bounded}
  \Longrightarrow\\
&\quad
  \Bigl[
    \limsup_{n\to\infty}x_n
    \text{ is a cluster value of }(x_n)
    \land
    \forall L\in\mathbb{R}\;
    \bigl(L\text{ is a cluster value of }(x_n)
    \Longrightarrow L\le \limsup_{n\to\infty}x_n\bigr)
  \Bigr]\\
&\quad\land
  \Bigl[
    \liminf_{n\to\infty}x_n
    \text{ is a cluster value of }(x_n)
    \land
    \forall L\in\mathbb{R}\;
    \bigl(L\text{ is a cluster value of }(x_n)
    \Longrightarrow \liminf_{n\to\infty}x_n\le L\bigr)
  \Bigr]
\Bigr).
\end{aligned}
\]
\end{remark*}

\begin{remark*}[Theorem predicate reading]
\[
\begin{aligned}
&\operatorname{BoundedSequence}(x_n)
\Longrightarrow\\
&\quad
\operatorname{SequenceClusterValue}
\bigl(x_n,\limsup_{n\to\infty}x_n\bigr)
\land
\forall L\in\mathbb{R}\;
\bigl(\operatorname{SequenceClusterValue}(x_n,L)
\Longrightarrow L\le \limsup_{n\to\infty}x_n\bigr),\\
&\operatorname{BoundedSequence}(x_n)
\Longrightarrow\\
&\quad
\operatorname{SequenceClusterValue}
\bigl(x_n,\liminf_{n\to\infty}x_n\bigr)
\land
\forall L\in\mathbb{R}\;
\bigl(\operatorname{SequenceClusterValue}(x_n,L)
\Longrightarrow \liminf_{n\to\infty}x_n\le L\bigr).
\end{aligned}
\]
\end{remark*}

\begin{remark*}[Interpretation]
Limsup and liminf describe the top and bottom edges of the sequence's
subsequential behavior.
\end{remark*}

\begin{remark*}[Dependencies]\hfill
\begin{itemize}
  \item \hyperref[def:limsup-sequence]{Limit Superior of a Sequence}
  \item \hyperref[def:liminf-sequence]{Limit Inferior of a Sequence}
  \item \hyperref[thm:limsup-largest-subsequential-limit]{Limsup Is the Largest Subsequential Limit}
  \item \hyperref[thm:liminf-smallest-subsequential-limit]{Liminf Is the Smallest Subsequential Limit}
  \item \hyperref[thm:cluster-values-are-subsequential-limits]{Cluster Values Are Subsequential Limits}
\end{itemize}
\end{remark*}